\subsection*{Physical scope of the icosahedral Standard Model and \(W/Z\) results}

The finite icosahedral package is an exact conditional recognition result.
Incidence determines the antipode \(J\). Under the explicit contract that an
admissible response is a signed central involutive graph automorphism
implementing inverse-port readback, the responses are exactly \(\pm J\);
their common sign is conventional. Conditional also on the rank-15 matter
contract with its unique charge-conjugate projector pair, the port-current,
anomaly, Spin, and tensor-descent certificates fix the \(3+2\) block
structure, hypercharge lattice, and \(\mathbb Z_6\) quotient from the stated
premises alone. The finite source packet therefore selects the odd-Weyl Spin
typing and the maximal faithful \(\mathbb Z_6\) quotient on its declared
carrier. The kernel is computed on every realized tensor and is insensitive
to the projector representative. This finite selection supplies no laboratory
identification of its current or flux sectors and no physical seam action.
Equality with the independently reconstructed Tannaka current, attachment of
the band to three physical chiral families, exclusion of extra light sectors,
four-dimensional topological attachment, scalar attachment and dynamics, and
construction of a chiral quantum field theory remain work in progress.

The conditional field-theory implications are explicit. A finite local action
gives an exact finite gauge-invariance and locality theorem, and the familiar
electroweak tree kernel is conditional on a separate canonical continuum and
action-normalization bridge. An exact finite measure criterion and an exact
finite Hamiltonian criterion are two parallel branches over that action. A
separate formal perturbative branch carries the strict finite-order \(W/Z\)
pole theorem. A nonperturbative continuum completion gives an observable-sector
reconstruction implication and a distinct continued-sheet resonance-stability
implication. The measure and perturbative branches are parallel descendants of
the finite local action. Neither implies the other, and a perturbative pole
does not imply the continuum completion.

These quantum field theory (QFT) theorems state what follows from typed action
and quantization packets. They do not show that the target-free source emits
those packets. The source-selected action and normalization, complete measure,
target-clean perturbative matching, physical current amplitudes, source law
and covariance, uncertainty enclosure, operational clock, and continuum tower
with its continued-sheet packet are work in progress. Two external Standard
Model calculation engines provide a bounded validation surface for the
perturbative matching protocol. They are validation inputs rather than OPH
source producers. No source-native dimensionless or physical-unit \(W/Z\)
pole follows from them.
